%==============================================================================
%  Lab Report Template — Lab 3: Unsupervised Learning
%  Course : Introduction to Machine Learning (with Python & scikit-learn)
%  Follows: notebooks/03-unsupervised-learning-code.ipynb
%           "Preprocessing, PCA & Clustering"
%
%  Compile: pdflatex lab3-unsupervised-learning.tex   (run twice for references)
%       or: latexmk -pdf lab3-unsupervised-learning.tex
%
%  HOW TO USE
%  ----------
%  1. Fill in the metadata in the "STUDENT INFORMATION" block below.
%  2. Run the notebook, then replace every  % TODO  / \textit{TODO} with the
%     numbers and plots you obtained.
%  3. Put your figures in a  figures/  subfolder and use \includegraphics.
%  4. Delete the guidance comments (lines starting with %) before submitting.
%==============================================================================
\documentclass[11pt,a4paper]{article}

%------------------------------ packages --------------------------------------
\usepackage[T1]{fontenc}
\usepackage[utf8]{inputenc}
\usepackage{lmodern}
\usepackage[margin=2.5cm]{geometry}
\usepackage{graphicx}
\usepackage{booktabs}
\usepackage{amsmath,amssymb}
\usepackage{xcolor}
\usepackage{listings}
\usepackage{caption}
\usepackage{subcaption}
\usepackage{enumitem}
\usepackage{float}
\usepackage{fancyhdr}
\usepackage[hidelinks]{hyperref}

%------------------------------ metadata --------------------------------------
% ========================= STUDENT INFORMATION ===============================
\newcommand{\Course}{Introduction to Machine Learning}
\newcommand{\LabNo}{Lab 3}
\newcommand{\LabTitle}{Unsupervised Learning}
\newcommand{\StudentName}{Your Name}          % TODO: your full name
\newcommand{\StudentID}{Student ID}           % TODO: your student number
\newcommand{\DueDate}{2026.10.21}                 % TODO: the due date
\newcommand{\Instructor}{Zhang Li \& Jing Ming}
\newcommand{\Institution}{Shandong Women's University}
% =============================================================================

%------------------------------ colours & code style --------------------------
\definecolor{codebg}{RGB}{245,242,247}
\definecolor{kw}{RGB}{125,15,98}      % keywords (burgundy)
\definecolor{str}{RGB}{31,51,82}      % strings (navy)
\definecolor{cmt}{RGB}{91,100,114}    % comments (muted)
\definecolor{accent}{RGB}{138,28,48}  % headings accent

\lstdefinestyle{py}{
  language=Python,
  backgroundcolor=\color{codebg},
  basicstyle=\ttfamily\small,
  keywordstyle=\color{kw}\bfseries,
  stringstyle=\color{str},
  commentstyle=\color{cmt}\itshape,
  numberstyle=\tiny\color{cmt},
  numbers=left,
  numbersep=8pt,
  frame=leftline,
  framerule=1.2pt,
  rulecolor=\color{accent},
  showstringspaces=false,
  breaklines=true,
  tabsize=4,
  xleftmargin=14pt,
  aboveskip=10pt,
  belowskip=6pt,
}
\lstset{style=py}

%------------------------------ headers/footers -------------------------------
\pagestyle{fancy}
\fancyhf{}
\lhead{\small \Course}
\rhead{\small \LabNo\ — \LabTitle}
\cfoot{\small \StudentName\ (\StudentID)}
\rfoot{\small \thepage}
\renewcommand{\headrulewidth}{0.4pt}
\renewcommand{\footrulewidth}{0.4pt}

%------------------------------ helpers ---------------------------------------
% A grey placeholder box to stand in for a figure you have not added yet.
\newcommand{\placeholder}[2][8cm]{%
  \fbox{\parbox[c][5cm][c]{#1}{\centering\color{cmt}\itshape #2}}%
}

%==============================================================================
\begin{document}

%------------------------------ title block -----------------------------------
\begin{titlepage}
\centering
\vspace*{2cm}
{\LARGE\bfseries \Course\par}
\vspace{0.6cm}
{\color{accent}\rule{\linewidth}{1.5pt}}\par
\vspace{0.5cm}
{\Huge\bfseries \LabNo: \LabTitle\par}
\vspace{0.3cm}
{\large Preprocessing, PCA \& Clustering on the Breast Cancer Dataset\par}
\vspace{0.4cm}
{\color{accent}\rule{\linewidth}{1.5pt}}\par
\vspace{2.3cm}

\begin{tabular}{rl}
  \textbf{Student:}    & \StudentName \\[4pt]
  \textbf{Student ID:} & \StudentID \\[4pt]
  \textbf{Instructor:} & \Instructor \\[4pt]
  \textbf{Institution:}& \Institution \\[4pt]
  \textbf{Date:}       & \DueDate \\
\end{tabular}

\vfill
{\color{cmt}\small Based on \texttt{notebooks/03-unsupervised-learning-code.ipynb}.
Replace all placeholder text before submission.}
\end{titlepage}

\tableofcontents
\newpage

%==============================================================================
\section{Objectives}
%------------------------------------------------------------------------------
In this lab you preprocess data and explore it \emph{without labels} using
dimensionality reduction and clustering. By the end you should be able to:
\begin{itemize}[leftmargin=*]
  \item \textbf{Scale} features correctly, fitting the scaler on the
        \textbf{training set only}.
  \item Reduce dimensions with \textbf{principal component analysis (PCA)} and
        read the \textbf{explained variance ratio}.
  \item Cluster data with \textbf{k-Means}, \textbf{agglomerative clustering},
        and \textbf{DBSCAN}.
  \item Evaluate clusterings against the true labels using the
        \textbf{adjusted Rand index (ARI)}.
\end{itemize}

%==============================================================================
\section{Background}
%------------------------------------------------------------------------------
In \emph{unsupervised learning} the algorithm receives only the inputs
$\mathbf{x}_i$ with \textbf{no target labels} and tries to discover structure in
the data. This lab covers two common tasks.

\paragraph{Feature scaling.} Many methods are sensitive to the \emph{scale} of
the features. \lstinline|StandardScaler| standardises each feature to zero mean
and unit variance,
\[
  z = \frac{x - \mu}{\sigma},
\]
where $\mu$ and $\sigma$ are the mean and standard deviation \textbf{estimated
on the training set}. Using test-set statistics would leak information, so the
scaler is \lstinline|fit| on the training data and then used to
\lstinline|transform| both sets.

\paragraph{PCA.} Principal component analysis finds orthogonal directions
(principal components) that capture the most variance in the data. Keeping the
first few components reduces dimensionality while retaining most information. The
\textbf{explained variance ratio} reports the fraction of total variance along
each component.

\paragraph{Clustering.} Clustering groups samples so that points in the same
group are more similar than points in different groups.
\begin{itemize}[leftmargin=*]
  \item \textbf{k-Means} partitions the data into $k$ clusters around centroids.
  \item \textbf{Agglomerative clustering} merges points bottom-up into a
        hierarchy until $k$ clusters remain.
  \item \textbf{DBSCAN} groups dense regions and marks low-density points as
        \textbf{noise} ($-1$); it chooses the number of clusters itself.
\end{itemize}

\paragraph{Evaluation.} Because we \emph{do} have the true labels here, we can
measure how well a clustering recovers them with the \textbf{adjusted Rand index}
(ARI). ARI ranges from $0$ (random labelling) to $1$ (perfect agreement); it is
adjusted so that chance groupings score about $0$.

%==============================================================================
\section{Dataset}
%------------------------------------------------------------------------------
We reuse the \textbf{Breast Cancer Wisconsin} dataset, bundled with scikit-learn
as \lstinline|load_breast_cancer()|. Each sample describes a tumour with 30
numeric features. The labels (\emph{malignant} / \emph{benign}) are used
\textbf{only for evaluation} (ARI) — the clustering algorithms never see them.

\begin{table}[H]
\centering
\caption{Dataset summary (fixed facts, provided for reference).}
\label{tab:dataset}
\begin{tabular}{@{}lr@{}}
\toprule
\textbf{Property} & \textbf{Value} \\
\midrule
Source              & \lstinline|sklearn.datasets.load_breast_cancer| \\
Task type           & Clustering / dimensionality reduction \\
Number of samples   & 569 \\
Number of features  & 30 (numeric) \\
Classes (eval only) & malignant (0), benign (1) \\
Split               & \lstinline|train_test_split(..., random_state=0)| \\
Training samples    & 426 \\
Test samples        & 143 \\
\bottomrule
\end{tabular}
\end{table}

%==============================================================================
\section{Methodology}
%------------------------------------------------------------------------------
The notebook is organised as five exercises. All tools come from scikit-learn.
The setup:

\begin{lstlisting}[caption={Imports and data (Setup).},label={lst:setup}]
import numpy as np
from sklearn.datasets import load_breast_cancer
from sklearn.model_selection import train_test_split
from sklearn.preprocessing import StandardScaler
from sklearn.decomposition import PCA
from sklearn.cluster import KMeans, AgglomerativeClustering, DBSCAN
from sklearn.metrics import adjusted_rand_score

cancer = load_breast_cancer()
\end{lstlisting}

\noindent The five exercises are:
\begin{enumerate}[leftmargin=*]
  \item \textbf{Scale correctly.} Split the data (\lstinline|random_state=0|),
        fit a \lstinline|StandardScaler| on the training set \emph{only}, and
        transform both sets into \lstinline|X_train_s|, \lstinline|X_test_s|.
  \item \textbf{PCA to 2 components.} Fit \lstinline|PCA(n_components=2)| on
        \lstinline|X_train_s|; report the explained variance ratio and the
        transformed shape.
  \item \textbf{k-Means.} Scale the whole dataset, run
        \lstinline|KMeans(n_clusters=2, n_init=10, random_state=0)|, and compute
        ARI against the labels.
  \item \textbf{Agglomerative clustering.} Run
        \lstinline|AgglomerativeClustering(n_clusters=2)| on the scaled data and
        compute its ARI.
  \item \textbf{DBSCAN.} Run \lstinline|DBSCAN()| on the scaled data; report the
        number of clusters found and the number of noise points ($-1$).
\end{enumerate}

%==============================================================================
\section{Experiments and Results}
%------------------------------------------------------------------------------
% Run the notebook and fill in every \textit{TODO} with the value you obtained.

\subsection{Exercise 1 — Scale correctly}
Fit the scaler on the training set only, then check that the scaled training
features have mean $\approx 0$ and standard deviation $\approx 1$.
\begin{itemize}[leftmargin=*]
  \item Scaled train mean (first 3 features): \textit{TODO} \quad(expected $\approx[0,0,0]$)
  \item Scaled train std (first 3 features): \textit{TODO} \quad(expected $\approx[1,1,1]$)
\end{itemize}

\subsection{Exercise 2 — PCA to 2 components}
Report the explained variance ratio and the shape after transforming.
\begin{table}[H]
\centering
\caption{PCA with two components (fitted on the scaled training set).}
\label{tab:pca}
\begin{tabular}{@{}lc@{}}
\toprule
\textbf{Quantity} & \textbf{Value} \\
\midrule
Explained variance ratio (PC1) & \textit{TODO} \\
Explained variance ratio (PC2) & \textit{TODO} \\
Total variance explained       & \textit{TODO} \\
Shape after PCA                & \textit{TODO} \quad(expected $(426, 2)$) \\
\bottomrule
\end{tabular}
\end{table}
% Optional: scatter plot of the two principal components, coloured by the label.
\begin{figure}[H]
\centering
% \includegraphics[width=0.7\linewidth]{figures/pca_scatter.png}
\placeholder[10cm]{(Optional) Scatter of PC1 vs.\ PC2,\\ coloured by the true label.}
\caption{Data projected onto the first two principal components (optional).}
\label{fig:pca}
\end{figure}

\subsection{Exercises 3 \& 4 — k-Means and agglomerative clustering}
Both algorithms are asked for \lstinline|n_clusters=2| on the fully scaled data.
Report the ARI of each against the true labels.
\begin{table}[H]
\centering
\caption{Clustering agreement with the true labels (adjusted Rand index).}
\label{tab:ari}
\begin{tabular}{@{}lc@{}}
\toprule
\textbf{Algorithm} & \textbf{ARI} \\
\midrule
\lstinline|KMeans(n_clusters=2, n_init=10, random_state=0)| & \textit{TODO} \\
\lstinline|AgglomerativeClustering(n_clusters=2)|           & \textit{TODO} \\
\bottomrule
\end{tabular}
\end{table}

\subsection{Exercise 5 — DBSCAN}
Run \lstinline|DBSCAN()| with its default parameters on the scaled data.
\begin{table}[H]
\centering
\caption{DBSCAN result with default parameters.}
\label{tab:dbscan}
\begin{tabular}{@{}lc@{}}
\toprule
\textbf{Quantity} & \textbf{Value} \\
\midrule
Number of clusters found & \textit{TODO} \\
Number of noise points ($-1$) & \textit{TODO} \\
\bottomrule
\end{tabular}
\end{table}

%==============================================================================
\section{Discussion}
%------------------------------------------------------------------------------
Answer with evidence from your tables and figures above:
\begin{itemize}[leftmargin=*]
  \item \textbf{Scaling (Ex.~1).} Why must the scaler be fitted on the training
        set only? What could go wrong if you fit it on the whole dataset?
  \item \textbf{PCA (Ex.~2).} How much of the total variance do the first two
        components capture? Is two dimensions enough to see the structure?
  \item \textbf{Clustering quality (Ex.~3--4).} Which algorithm's clusters agreed
        best with the true labels? Are k-Means and agglomerative clustering close?
  \item \textbf{DBSCAN (Ex.~5).} With default parameters, how many clusters did
        DBSCAN find and how many points were labelled noise? What does this tell
        you about choosing \lstinline|eps| and \lstinline|min_samples|?
  \item \textbf{Labels vs.\ structure.} The clusters never saw the labels. What
        does a high ARI suggest about the geometry of the two tumour classes?
\end{itemize}
% TODO: write 1--2 paragraphs answering the questions above.

%==============================================================================
\section{Conclusion}
%------------------------------------------------------------------------------
% 3--5 sentences summarising what you learned and the main result.
% TODO.

%==============================================================================
% \begin{thebibliography}{9}
% \bibitem{mueller2016}
% A.~C. M\"uller and S.~Guido,
% \emph{Introduction to Machine Learning with Python}.
% O'Reilly Media, 2016.

% \bibitem{sklearn}
% F.~Pedregosa et al.,
% ``Scikit-learn: Machine Learning in Python,''
% \emph{Journal of Machine Learning Research}, vol.~12, pp.~2825--2830, 2011.
% \end{thebibliography}

%==============================================================================
\appendix
\section{Source Code}
%------------------------------------------------------------------------------
% Paste your completed code, or reference the notebook you ran.
Full code: \texttt{notebooks/03-unsupervised-learning-code.ipynb}. Key solution:

\begin{lstlisting}[caption={Reference solution for Exercises 1--5.}]
X_train, X_test, y_train, y_test = train_test_split(
    cancer.data, cancer.target, random_state=0)

# Ex1 - scale, fitting on the training set only
scaler = StandardScaler().fit(X_train)
X_train_s = scaler.transform(X_train)
X_test_s  = scaler.transform(X_test)

# Ex2 - PCA to 2 components
pca = PCA(n_components=2).fit(X_train_s)
print("explained variance ratio:", np.round(pca.explained_variance_ratio_, 3))
print("shape after PCA:", pca.transform(X_train_s).shape)

# Ex3-4 - k-Means and agglomerative clustering on the whole scaled dataset
X_scaled = StandardScaler().fit_transform(cancer.data)
km  = KMeans(n_clusters=2, n_init=10, random_state=0).fit(X_scaled)
ari_kmeans = adjusted_rand_score(cancer.target, km.labels_)
agg = AgglomerativeClustering(n_clusters=2).fit(X_scaled)
ari_agg = adjusted_rand_score(cancer.target, agg.labels_)
print("ARI k-means      :", round(ari_kmeans, 3))
print("ARI agglomerative:", round(ari_agg, 3))

# Ex5 - DBSCAN
db = DBSCAN().fit(X_scaled)
labels = db.labels_
n_clusters = len(set(labels)) - (1 if -1 in labels else 0)
print("DBSCAN clusters:", n_clusters,
      "| noise points:", int((labels == -1).sum()))
\end{lstlisting}

\end{document}
